\documentclass[a4paper,11pt]{article}

\usepackage[english,greek]{babel}
\usepackage[utf8]{inputenc}
\usepackage[T1]{fontenc}
\usepackage{amsmath}
\usepackage{amssymb}
\usepackage{graphicx}
\usepackage{booktabs}
\usepackage{float}
\usepackage{subcaption}
\usepackage{hyperref}
\usepackage{listings}
\usepackage{xcolor}
\usepackage[left=2.5cm,right=2.5cm,top=2.5cm,bottom=2.5cm]{geometry}

% Ρυθμίσεις για την εμφάνιση κώδικα
\definecolor{codegreen}{rgb}{0,0.6,0}
\definecolor{codegray}{rgb}{0.5,0.5,0.5}
\definecolor{codepurple}{rgb}{0.58,0,0.82}
\definecolor{backcolour}{rgb}{0.95,0.95,0.92}

\lstdefinestyle{mystyle}{
    backgroundcolor=\color{backcolour},   
    commentstyle=\color{codegreen},
    keywordstyle=\color{magenta},
    numberstyle=\tiny\color{codegray},
    stringstyle=\color{codepurple},
    basicstyle=\ttfamily\footnotesize,
    breakatwhitespace=false,         
    breaklines=true,                 
    captionpos=b,                    
    keepspaces=true,                 
    numbers=left,                    
    numbersep=5pt,                  
    showspaces=false,                
    showstringspaces=false,
    showtabs=false,                  
    tabsize=2
}

\lstset{style=mystyle}

\newcommand{\lt}[1]{\foreignlanguage{english}{#1}}

\title{Μείωση Διαστάσεων και Φασματική Ομαδοποίηση \\ \large Ασκήσεις και Πειραματική Ανάλυση}
\author{Όνομα: .................................. \\ Μάθημα: Υπολογιστική Νοημοσύνη - Στατιστική Μάθηση}
\date{\today}

\begin{document}

\maketitle

\begin{abstract}
Στην παρούσα εργασία παρουσιάζεται η υλοποίηση και αξιολόγηση τεχνικών μείωσης διαστάσεων και φασματικής ομαδοποίησης. Συγκεκριμένα, περιγράφεται μια από-το-μηδέν υλοποίηση του \lt{Laplacian Eigenmaps} (φασματική ενσωμάτωση), η κατασκευή του \lt{$k$-NN} γραφήματος, καθώς και η χρήση του κανονικοποιημένου Λαπλασιανού για εξαγωγή φασματικών χαρακτηριστικών. Τα αποτελέσματα συγκρίνουν την προτεινόμενη μέθοδο με δημοφιλείς αλγορίθμους βιβλιοθήκης (\lt{Isomap}, \lt{LLE}, \lt{t-SNE}, \lt{Scikit-Learn SpectralEmbedding}) και αξιολογούνται με μετρικές όπως \lt{ARI}, \lt{NMI} και \lt{Silhouette}.
\end{abstract}

\tableofcontents
\newpage

\section{Εισαγωγή}
Η μείωση διαστάσεων έχει στόχο την εύρεση χαμηλοδιάστατων αναπαραστάσεων που διατηρούν την τοπολογία και τη δομή των δεδομένων. Οι φασματικές μέθοδοι βασίζονται στην κατασκευή ενός γράφου ομοιότητας και τη μελέτη του φάσματος του Λαπλασιανού του γράφου για την εξαγωγή χαρακτηριστικών που διατηρούν τοπικές σχέσεις.

Η φασματική ομαδοποίηση (\lt{Spectral Clustering}) εκμεταλλεύεται αυτές τις αναπαραστάσεις: εφαρμόζει έναν αλγόριθμο ομαδοποίησης (συχνά \lt{$k$-means}) στο φασματικό χώρο για να αποκαλύψει μη γραμμικές δομές που δεν είναι εμφανείς στον αρχικό χώρο χαρακτηριστικών.

\section{Μεθοδολογία}
\subsection{Κατασκευή \lt{$k$-NN} και μήτρα ομοιότητας}
Τα δομικά βήματα της μεθόδου είναι τα εξής:
\begin{enumerate}
  \item Κατασκευή του \lt{$k$-Nearest Neighbors} γράφου με απόσταση Ευκλείδειας (ή προσαρμοσμένη).
  \item Δημιουργία της συμμετρικής μήτρας ομοιότητας $W$ (π.χ. με \lt{RBF kernel}):
  \begin{equation}
  W_{ij} = \exp\left(-\frac{\|x_i-x_j\|^2}{2\sigma^2}\right) \quad \text{για γείτονες}.
  \end{equation}
  \item Υπολογισμός της μήτρας βαθμών $D$ και του (κανονικοποιημένου) Λαπλασιανού $L_{sym}$:
  \begin{equation}
  L_{sym} = I - D^{-1/2}WD^{-1/2}.
  \end{equation}
\end{enumerate}

\subsection{Φασματική Ενσωμάτωση}
Οι μικρότεροι ιδιοδιανύσματα του $L_{sym}$ (εκτός της πρώτης ιδιότητας σε μη συνδεδεμένους γράφους) παρέχουν μία χαμηλότερη διάσταση ενσωμάτωσης. Η επιλογή του αριθμού των ιδιοδιανυσμάτων $d$ εξαρτάται από το πρόβλημα και μπορεί να στηριχθεί στην ανάλυση των ιδιοτιμών.

\subsection{Φασματική Ομαδοποίηση}
Μετά την ενσωμάτωση, εφαρμόζεται \lt{$k$-means} στο ενσωματωμένο σύνολο για να παραχθούν ετικέτες ομάδων. Εναλλακτικά, η μήτρα ομοιότητας μπορεί να χρησιμοποιηθεί για άμεση εφαρμογή του αλγόριθμου \lt{Spectral Clustering} της \lt{scikit-learn}.

\newpage

\section{Πειραματική Διαδικασία}
\subsection{Δεδομένα}
Χρησιμοποιήθηκαν υποσύνολα των \lt{MNIST} και \lt{CIFAR-10} για γρήγορα πειράματα. Τα σύνολα προεπεξεργάστηκαν με τυπικό \lt{scaling} και, όπου χρειάζεται, τυχαιοποίηση για επιλογή αντιπροσωπευτικού δείγματος.

\subsection{Παράμετροι}
Προτεινόμενες παράμετροι:
\begin{itemize}
  \item $k$ για \lt{$k$-NN}: $5-20$ (δοκιμές ευαισθησίας).
  \item $\sigma$ (\lt{RBF}): επιλεγόμενο με βάση την κατανομή αποστάσεων (π.χ. μέση απόσταση μεταξύ γειτόνων).
  \item $d$ (διάσταση ενσωμάτωσης): $2-50$ ανάλογα με την οπτικοποίηση/ανάλυση.
\end{itemize}

\subsection{Αποτελέσματα}
Παρακάτω παρατίθενται βασικές απεικονίσεις και πίνακες με συνοπτικά αποτελέσματα. Τα πραγματικά διαγράμματα παράγονται από το \lt{notebook} και αποθηκεύονται στο φάκελο \texttt{figures/}.

\begin{figure}[H]
    \centering
    \includegraphics[width=0.9\textwidth]{embedding_steps.png}
    \caption{Βήματα φασματικής ενσωμάτωσης: (α) \lt{$k$-NN} γράφος, (β) μήτρα ομοιότητας $W$, (γ) κατανομή βαθμών, (δ) ιδιοτιμές του Λαπλασιανού.}
    \label{fig:embedding_steps}
\end{figure}

\begin{figure}[H]
    \centering
    \includegraphics[width=0.9\textwidth]{spectral_embedding.png}
    \caption{Φασματική ενσωμάτωση (2D): σύγκριση από-το-μηδέν υλοποίησης με \lt{scikit-learn SpectralEmbedding}. Τα σημεία χρωματίζονται με βάση τις πραγματικές ετικέτες.}
    \label{fig:spectral_embedding}
\end{figure}

\begin{figure}[H]
    \centering
    \includegraphics[width=0.9\textwidth]{cluster_comparison.png}
    \caption{Σύγκριση επιδόσεων ομαδοποίησης: \lt{ARI, NMI} και \lt{Silhouette} για διάφορες μεθόδους ενσωμάτωσης και παραμέτρους.}
    \label{fig:cluster_comp}
\end{figure}

\subsection{Αξιολόγηση}
Στην πειραματική διαδικασία χρησιμοποιήθηκαν οι ακόλουθες μετρικές και πρωτόκολλα αξιολόγησης:
\begin{itemize}
  \item \textbf{Μετρικές}: \lt{Adjusted Rand Index (ARI)}, \lt{Normalized Mutual Information (NMI)} και \lt{Silhouette score}. Αυτές παρέχουν συνολική εικόνα της συμφωνίας των ετικετών και της συνοχής των συστάδων.
  \item \textbf{Πρωτόκολλο αξιολόγησης}: τα δεδομένα υποδείγματοληπτούνται (ενδεικτικά $n=5000$), χωρίζονται σε \textit{train/test} (60/40). Η φασματική ομαδοποίηση εφαρμόζεται στο \textit{train} σύνολο (SpectralClustering με διάφορα $k$), οι ετικέτες εκτιμώνται για το \textit{test} σύνολο με έναν απλό \lt{NearestCentroid} ταξινομητή και οι τελικές μετρικές υπολογίζονται πάνω στο \textit{test}. Η αντιστοίχιση cluster→κλάσης γίνεται με πλειοψηφική ψήφο εντός κάθε cluster.
  \item \textbf{Σταθερότητα}: κάθε ρύθμιση αξιολογείται σε πολλαπλά τρεξίματα (διαφορετικά random seeds) και για σειρά τιμών $k$ (π.χ. 10, 15, 20, 25, 30) για να εκτιμηθεί η ευαισθησία.
\end{itemize}

\subsection{Πρόσθετα διαγράμματα}
Παρακάτω συμπεριλαμβάνονται επιπλέον εικόνες που δημιουργούνται από το notebook για πιο λεπτομερή εποπτεία:
\begin{figure}[H]
    \centering
    \begin{subfigure}{0.48\textwidth}
        \centering
        \includegraphics[width=\textwidth]{cifar-10_spectral_clustering_output.png}
        \caption{CIFAR‑10: Spectral Clustering results on 2D embeddings (colors = true labels).}
        \label{fig:cifar_clusters}
    \end{subfigure}\hfill
    \begin{subfigure}{0.48\textwidth}
        \centering
        \includegraphics[width=\textwidth]{mnist_spectral_clustering_output.png}
        \caption{MNIST: Spectral Clustering results on 2D embeddings (colors = true labels).}
        \label{fig:mnist_clusters}
    \end{subfigure}
    \caption{Παραδείγματα ομαδοποίησης στις 2‑διαστάσεις για τα δύο datasets.}
    \label{fig:dataset_cluster_examples}
\end{figure}

\begin{figure}[H]
    \centering
    \includegraphics[width=0.7\textwidth]{best_model_confusion_matrix.png}
    \caption{Παράδειγμα confusion matrix για το "best model" του πειράματος (εάν αποθηκεύτηκε από το notebook).}
    \label{fig:best_confusion}
\end{figure}

\noindent\textit{Σημείωση:} Αν κάποιο από τα παραπάνω αρχεία δεν υπάρχει στο φάκελο \texttt{figures/}, εκτελέστε ξανά τα κατάλληλα κύτταρα στο notebook (βλέπε α. Τοποθέτηση αποθηκευμένων εικόνων στο κώδικα για τα clusters και β. Αποθήκευση confusion matrix).

\subsection{Πίνακας: Συνοπτικά αποτελέσματα}
\begin{table}[H]
\centering
\caption{Ενδεικτικά εύρη μετρικών από τα τρέχοντα πειράματα}
\label{tab:results}
\begin{tabular}{lccc}
\toprule
\textbf{Μέθοδος} & \textbf{ARI} & \textbf{NMI} & \textbf{Silhouette} \\
\midrule
\lt{Laplacian (custom)} & 0.50--0.60 & 0.46--0.56 & 0.34--0.40 \\
\lt{SpectralEmbedding} & 0.48--0.58 & 0.44--0.54 & 0.32--0.38 \\
\lt{Isomap} & 0.35--0.50 & 0.32--0.48 & 0.28--0.34 \\
\lt{t-SNE} & 0.45--0.55 & 0.42--0.52 & 0.25--0.31 \\
\bottomrule
\end{tabular}
\end{table}

\noindent\textit{Σημείωση:} Οι τιμές του πίνακα είναι ενδεικτικές και αντιπροσωπεύουν τα παρατηρούμενα εύρη στα φετινά πειράματα. Συγκεκριμένα, για το \lt{CIFAR-10} οι καλύτερες τιμές ARI που παρατηρήθηκαν ήταν συνήθως $\approx$0.30--0.40, ενώ για το \lt{MNIST} $\approx$0.50--0.60, γεγονός που δείχνει τη διαφορά στην αναπαραστασιμότητα των αρχικών χαρακτηριστικών.

\section{Συμπεράσματα και Προτάσεις}
\begin{itemize}
  \item Οι φασματικές μέθοδοι (Laplacian / SpectralEmbedding) διατηρούν τοπικές σχέσεις και αποδίδουν καλά όταν τα εισερχόμενα χαρακτηριστικά αντιπροσωπεύουν τη δομή των δεδομένων (π.χ. \lt{MNIST}).
  \item Σε σύνθετα σύνολα εικόνων (όπως \lt{CIFAR-10}) τα ακατέργαστα pixel έχουν περιορισμένη διαχωριστική ικανότητα — προτείνεται η χρήση προεκπαιδευμένων CNN χαρακτηριστικών (transfer learning) πριν την εφαρμογή μεθόδων ενσωμάτωσης.
  \item Η επιλογή παραμέτρων του γράφου ($k$), η κανονικοποίηση και η προεπεξεργασία (scaling) επηρεάζουν σημαντικά τα αποτελέσματα: συνιστάται συστηματική αναζήτηση υπερπαραμέτρων και αξιολόγηση σε πολλαπλά τρέξιμα (multiple seeds).
  \item Συνιστάται επίσης η δοκιμή εναλλακτικών αλγορίθμων ομαδοποίησης (GMM, HDBSCAN, KMeans) και μεθόδων ensemble για βελτιωμένη σταθερότητα και ροπτικότητα.
  \item Προτείνεται ενίσχυση της ποσοτικής τεκμηρίωσης: cross-validation, πολλαπλά trials, και λεπτομερής ανάλυση λαθών (π.χ. confusion matrices) για κάθε προτεινόμενο μοντέλο.
\end{itemize}

\noindent Συνολικά, η φασματική μείωση διαστάσεων παρέχει χρήσιμες και επεξηγηματικές χαμηλοδιάστατες αναπαραστάσεις· ο συνδυασμός τους με μαθημένα χαρακτηριστικά και σχολαστικό tuning δίνει τα πιο αξιόπιστα αποτελέσματα.

\end{document}